\section{Local criteria of flatness}
\label{IV.5}

\begin{proposition}
\label{IV.5.1}
Let $A$ be a ring equipped with an ideal $I$, $M$ an
$A$-module. Suppose
$$
\Tor_1^A(M, A/I^n) = 0 \qquad \text{for $n>0$}
$$
then the canonical surjective homomorphism
\begin{equation*}
\label{eq:IV.5.1.*}
\tag{$*$} {\gr_I^0(M) \otimes_{A/I} \gr_I(A) \to \gr_I(M)}
\end{equation*}
is an isomorphism. The converse is true if $I$ is nilpotent.
\end{proposition}

The hypothesis means that the homomorphisms
$$
I^n \otimes_A M \to I^n M
$$
are isomorphisms, whence at once the fact that the
homomorphisms
$$
I^n / I^{n+1} \otimes_A M \to I^n M / I^{n+1} M
$$
are isomorphisms. Conversely, suppose this condition
satisfied and $I$ nilpotent; let us prove $\Tor_1^A(M,A/I^n) =0$
for every $n$. It is true for $n$ large; let us proceed by
descending induction on~$n$, supposing it proved for
$n+1$. We have a commutative diagram
$$
\xymatrix{ & M \otimes I^{n+1} \ar[r] \ar[d] & M \otimes I^n \ar[r]
\ar[d] & M \otimes(I^n/I^{n+1}) \ar[r] \ar[d] & 0 \\ 0 \ar[r] &
MI^{n+1} \ar[r] & MI^n \ar[r] & MI^n/MI^{n+1} \ar[r] & 0 }
$$
where
% original p. 96
the rows are exact. By hypothesis, the last vertical
arrow is an isomorphism, and the induction hypothesis
also means that the first vertical arrow is one. Hence the
same holds for the middle vertical arrow, which
completes the proof.

The following proposition was brought out at the time of the
Seminar by Serre; it allows substantial simplifications
in the present no.
\begin{proposition}
\label{IV.5.2}
Let $A \to B$ be a ring homomorphism, $M$ an $A$-module. The
following conditions are equivalent:
\begin{enumerate}
\item[(i)] For every $B$-module $N$, one has $\Tor_1^A(M,N) = 0$;
\item[(ii)] $\Tor_1^A(M,B) = 0$, and $M_{(B)} = M \otimes_A B$ is
$B$-flat.
\end{enumerate}
\end{proposition}

We have a functorial isomorphism
$$
M \otimes_A N = (M \otimes_A B) \otimes_B N
$$
which expresses the left-hand side, considered as a functor in
$M$, as a composite of two functors $M \mto M \otimes_A B$
and $P \mto P \otimes_B N$. As the first transforms free
modules over~$A$ into free modules over~$B$, hence projectives into
projectives, we have the spectral sequence of composite functors
$$
\Tor_n^A(M,N) \From \Tor_p^B(\Tor_q^A(M,B),N)
$$
whence an exact sequence for the terms of low degree
$$
0 \from \Tor_1^B(M \otimes_A B,N) \from \Tor_1^A(M,N)
\from \Tor_1^A(M,B) \otimes_A N
$$
If (i) is satisfied, then one concludes from this exact sequence
$\Tor_1^B(M \otimes_A B,N) = 0$ for every $N$, \ie $M \otimes_A B$
is $B$-flat, whence (ii). If conversely (ii) is satisfied,
then in the exact sequence the terms surrounding $\Tor_1^A(M,N)$ are
zero, hence we have (i).

\begin{corollary}
\label{IV.5.3}
Suppose that $B = A/I$; then the preceding conditions
are equivalent to the following one:
\begin{enumerate}
\item[(iii)] $\Tor_1^A(M,N) = 0$ for every $A$-module $N$ annihilated
by a power of~$I$.
\end{enumerate}
\end{corollary}

Indeed, (i) means that this is so if $N$ is annihilated by
$I$. One deduces (iii) from this by applying the hypothesis to the
$I^nN/I^{n+1}N$.

\begin{corollary}
\label{IV.5.4}
Under the conditions of~\Ref{IV.5.3}, the conditions considered
imply that
% original p. 97
the functorial ho\-mo\-mor\-phism
\begin{equation*}
\label{eq:IV.5.*}
\tag{$*$} {\gr_I^0(M) \otimes_{A/I} \gr_I(A) \to \gr_I(M)}
\end{equation*}
is an isomorphism, and that $M \otimes_A A/I$ is flat over~$A/I$.
\end{corollary}

It suffices to apply (iii) and \eqref{IV.5.1}. Using the
converse of~\Ref{IV.5.1} in the case $I$ nilpotent, one finds:
\begin{corollary}
\label{IV.5.5}
Let $A$ be a ring equipped with a nilpotent ideal $I$, $M$ an
$A$-module. The following conditions are equivalent:
\begin{enumerate}
\item[(i)] $M$ is $A$-flat
\item[(ii)] $M \otimes_A A/I$ is $A/I$-flat, and $\Tor_1^A(M,A/I) = 0$
\item[(iii)] $M \otimes_A A/I$ is $A/I$-flat, and the canonical
homomorphism \eqref{eq:IV.5.*} on the graded modules is an isomorphism.
\end{enumerate}
\end{corollary}

Indeed, these are respectively the preceding conditions (iii) and (ii),
and those of corollary~\Ref{IV.5.4}.

Let us no longer suppose $I$ nilpotent; then we shall only have a priori
in~\Ref{IV.5.5} the implications (i)$\To$(ii)$\To$(iii). On the other hand, as condition (iii) remains
stable on dividing by a power of~$I$, we see by virtue
of~\Ref{IV.5.5} that it implies
\begin{enumerate}
\item[(iv)] \emph{for every integer $n$, $M \otimes A/I^n$ is flat over~$A/I^n$.}
\end{enumerate}
We propose to give conditions under which one can
conclude from this~(i), \ie that $M$ is $A$-flat. I say that it suffices for
this that $A$ be noetherian and that~$M$ satisfy the following
finiteness condition: \emph{for every module of finite type $N$ over~$A$, $M
\otimes_A N$ is separated for the
$I$-preadic topology}. (It would suffice to verify it if~$N$ is an
ideal of finite type in $A$). Indeed, let us prove that under these
conditions, if $N' \to N$ is a monomorphism of modules of finite
type, $M \otimes_A N' \to M \otimes_A N$ is a monomorphism. It
suffices indeed to show that the kernel is contained in the $I^n(M
\otimes_A N') = \Im(M \otimes_A I^n N' \to M \otimes_A N')$, or again
in the $\Im(M \otimes_A V'_n \to M \otimes_A N') = \Ker (M \otimes_A
N' \to M \otimes_A (N'/V'_n))$, where $V'_n$ runs through a countable
fundamental system of neighborhoods of 0 in $N'$ (equipped with its
$I$-adic topology). By Krull's theorem, the
$I$-adic topology of~$N'$ is induced by that of~$N$, so we can
take $V'_n = N' \cap I^n N$. Consider then the
commutative diagram
$$
\xymatrix{ M \otimes_A N' \ar[r] \ar[d] & M \otimes_A (N'/V'_n) \ar[d]
\\ M \otimes_A N \ar[r] & M \otimes_A (N / I^n N) }
$$
% original p. 98
As $N'/V'_n$ and $N/I^n N$ are annihilated by $I^n$, the
second vertical homomorphism is identified with the one deduced
from the \emph{injective} homomorphism $N'/V'_n \to N/I^n N$ by
tensoring over~$A/I^n$ with the \emph{flat} $(A/I^n)$-module $M
\otimes_A A/I^n$; it is therefore \emph{injective}. Consequently, the kernel of
$M \otimes_A N' \to M \otimes_A N$ is contained in that of~$M
\otimes_A N' \to M \otimes_A (N'/V'_n),$ which is what we wanted.

The ``finiteness'' condition considered on~$M$ is
satisfied in particular if $M$ is a module of finite type over
a noetherian $A$-algebra
$B$ such that $IB$ is contained
in the radical of~$B$: indeed, then $M \otimes_A N$ is a module
of finite type over~$B$ for every module of finite type $N$ over~$A$, hence
separated by Krull for the $I$-adic topology $=$ its
$(IB)$-adic topology. We thus find:

\begin{theorem}
\label{IV.5.6}
Let $A\to B$ be a homomorphism of noetherian rings, $I$ an
ideal of~$A$ such that $IB$ is contained in the radical of~$B$, $M$
a $B$-module of finite type. The following conditions are
equivalent:
%
\begin{enumerate}
\item[(i)] $M$ is $A$-flat
\item[(ii)] $M\otimes_A A/I$ is $A/I$-flat, and $\Tor_1^A(M,A/I)=0$
\item[(iii)] $M\otimes_A A/I$ is $A/I$-flat, and the canonical
homomorphism
%
$$
\gr_I^0(M)\otimes_{A/I}\gr_I(A)\to\gr_I(M)
$$
%
is an isomorphism.
\item[(iv)] For every integer $n$, $M\otimes_A A/I^n$ is flat
over~$A/I^n$.
\end{enumerate}
\end{theorem}

This result applies above all when $A$, $B$ are noetherian
\emph{local} rings, $A\to B$ a local homomorphism, and~$I$
an ideal of~$A$ contained in its maximal ideal (and one can
immediately reduce~\Ref{IV.5.6} to this case). An interesting
case is that where $A/I$ is a field, \ie $I$ maximal,
in which case the condition that $M\otimes_A(A/I)$ be flat over~$A/I$
becomes unnecessary; moreover, as then the $A/I^n$ are artinian
local rings, condition (iv) means that the
$M\otimes_A(A/I^n)$ are \emph{free} over the~$A/I^n$.

\begin{corollary}
\label{IV.5.7}
Let
% original p. 99
$A\to B$ be a local homomorphism of noetherian local rings,
$u\colon M'\to M$ a homomorphism of~$B$-modules of finite type;
suppose $M$ flat over~$A$. Then the following conditions are
equivalent:
%
\begin{enumerate}
\item[(i)] $u$ is injective, and $\Coker u$ is flat over~$A$.
\item[(ii)] $u\otimes_A k\colon M'\otimes_A k\to M\otimes_A k$ is
injective
\end{enumerate}
%
(where $k$ denotes the residue field of~$A$).
\end{corollary}

(i)$\To$(ii) by virtue of~\Ref{IV.1.1}; let us prove the
converse. First of all $u$ is injective, for it suffices to
verify this on the associated graded modules, where it follows from
a commutative square that the reader will write down. Let $M''$ be its
$\Coker$; we thus have an exact sequence
%
$$
0\to M'\to M\to M''\to 0
$$
%
whence by the exact sequence of the $\Tor$, taking into account
hypothesis~(ii) and~$\Tor_1^A(M,k)=0$, the relation
$\Tor_1^A(M'',k)=0$, hence $M''$ is flat over~$A$ by
theorem~\Ref{IV.5.6}.

\begin{corollary}
\label{IV.5.8}
Under the conditions of~\Ref{IV.5.6}, let $J$ be an ideal of~$B$
containing $IB$ and contained in the radical. Let $\hat A$ be the
$I$-adic completion of~$A$ and~$\hat B$ and~$\hat M$ the
$J$-adic completions of~$B$ and~$M$. For $M$ to be
$A$-flat, it is necessary and sufficient that $M$ be $\hat A$-flat.
\end{corollary}

(N.B. the sufficiency would already follow easily
from~\Ref{IV.3.2}). We use criterion (iii) of~\Ref{IV.5.6}
in the situation $(A,B,I,M)$ and in the situation $(\hat A,\hat
B,I\hat A, \hat M)$. One observes that the conditions obtained for
the one and the other case are equivalent, thanks
to~\Ref{IV.3.2}.

\begin{corollary}
\label{IV.5.9}
Let $A\to B\to C$ be local homomorphisms of noetherian local
rings, $M$ a $C$-module of finite type (N.B. $C$
intervenes only in order to be able to put a finiteness condition on~$M$). We suppose $B$ flat over~$A$. Let $k$ be the residue field
of~$A$. Equivalent conditions:
%
\begin{enumerate}
\item[(i)] $M$ is flat over~$B$.
\item[(ii)] $M$ is flat over~$A$, and $M\otimes_A k$ is flat over~$B\otimes_A k$.
\end{enumerate}
\end{corollary}

The implication (i)$\To$(ii) is trivial; let us prove
(ii)$\To$(i). We apply criterion~(iii)
of~\Ref{IV.5.6} to $(B,C,\goth{m}B=I,M)$; as
$M\otimes_B(B/I)=M\otimes_B(B\otimes_A k) = M\otimes_A k$, the
first condition of this criterion means precisely
that $M\otimes_A k$ is flat over~$B\otimes_A k$,
% original p. 100
fine. The second condition of the criterion is
satisfied because $M$ is flat over~$A$ and~$B$ flat over~$A$,
by an associativity formula for the tensor product. --- Of
course, referring to~\Ref{IV.5.5} instead
of~\Ref{IV.5.6}, one obtains an analogous statement without noetherian
and finiteness conditions, when one supposes on the other hand that
the ideal~$\goth{m}$ of~$A$ is nilpotent. (The fact that $\goth{m}$
was taken maximal has moreover not intervened; but
it is in a sense the case ``$\goth{m}$ maximal'' which is ``the best
possible'').

\section{Flat morphisms and open sets}
\label{IV.6}

Let us first recall some results on constructible sets,
which are moreover proved in notes
in circulation of the Dieudonn\'e-Rosenlicht Seminar on
Schemes\footnote{\Cf EGA~0$_\mathrm{III}$~9, EGA~IV~1.8 and~1.10.}.

Let $X$ be a topological space. We say with Chevalley that a subset
of~$X$ is \emph{constructible}
if it is a finite union of locally closed subsets.

\begin{lemma}
\label{IV.6.1}
Let $X$ be a noetherian topological space, let $Z$ be a subset
of~$X$. For~$Z$ to be constructible, it is necessary and sufficient that
for every irreducible closed subset $Y$ of~$X$, $Z\cap Y$
be non-dense in $Y$ or contain a nonempty open subset of
the space~$Y$.
\end{lemma}

One deduces from it, using a well-known lemma of Commutative
Algebra:

\begin{lemma}[Chevalley]
\label{IV.6.2}
Let $f\colon X\to Y$ be a morphism of finite type of
preschemes, with $Y$ noetherian. Then $f(X)$ is
constructible.
\end{lemma}

\begin{lemma}
\label{IV.6.3}
Let $X$ be a noetherian topological space every irreducible closed
subset of which admits a generic point, $U$ a
constructible subset of~$X$, $x\in X$. For $U$ to be a neighborhood
of~$x$, it is necessary and sufficient that every generization $y$ of~$x$
(\ie every $y\in X$ such that $x\in \bar{y}$) be in~$U$.
\end{lemma}

In particular

\begin{corollary}
\label{IV.6.4}
Let $X$ be a noetherian topological space every irreducible closed
subset of which admits a generic point, $U$ a
subset of~$X$. For $U$ to be open, it is necessary and sufficient
that it satisfy the following two conditions:
%
\begin{enumerate}
\item[(a)] $U$ contains every generization of each of its
points
\item[(b)] if $x\in U$, then
% original p. 101
$U\cap\bar{x}$ contains a nonempty open subset of
the space~$\bar{x}$.
\end{enumerate}
\end{corollary}

Indeed, $U$ is necessarily constructible thanks
to~\Ref{IV.6.1}, and we apply criterion~\Ref{IV.6.2}, which
proves that $U$ is a neighborhood of each of its points.

\begin{corollary}
\label{IV.6.5}
Let $f\colon X\to Y$ be a morphism of finite type of preschemes,
with $Y$ locally noetherian, $x$ a point of~$X$, $y=f(x)$.
For $f$ to transform every neighborhood of~$x$ into a neighborhood of~$y$,
it is necessary and sufficient that for every generization $y'$ of~$y$,
there exist a generization $x'$ of~$x$ such that $f(x')=y'$.
\end{corollary}

We can evidently suppose that $X$ and~$Y$ are affine, hence
noetherian. The condition is sufficient, for it suffices to
prove that $f(X)$ is a neighborhood of~$y$; now $f(X)$ is
constructible by~\Ref{IV.6.1}, and it suffices to apply
criterion~\Ref{IV.6.3}. The condition is necessary, for let
$Y'=\overline{y'}$, and let $F$ be the union of the irreducible
components of~$f^{-1}(Y')$ which do not contain $x$. Then
$X-F$ is an open neighborhood of~$x$, hence its image is a neighborhood
of~$y$, and a fortiori contains $y'$, hence there exists $x'_1\in X-F$ such
that $f(x'_1)=y'$. Consider an irreducible component
of~$f^{-1}(Y')$ containing $x'_1$; it necessarily contains $x$
(for otherwise it would be contained in $F$); let $x'$ be its
generic point. It is a generization of~$x$, and $f(x')$
is a generization of~$f(x'_1)=y'$ contained in~$Y'$, hence
is equal to $y'$, qed.

\begin{theorem}
\label{IV.6.6}
Let $f\colon X\to Y$ be a morphism locally of finite type, with~$Y$
locally noetherian, $F$ a coherent sheaf on~$X$ with
support~$X$, flat with respect to~$Y$. Then $f$ is an
open morphism (\iev transforms open sets into open sets).
\end{theorem}

It suffices to prove criterion~\Ref{IV.6.5} for every point $x\in
X$. Now the generizations $x'$ of~$x$ correspond to the
prime ideals of~$\cal{O}_x$, those~$y'$ of~$y$ correspond
to the prime ideals of~$\cal{O}_y$, and one must therefore verify
that every prime ideal of~$\cal{O}_y$ is induced by a prime
ideal of~$\cal{O}_x$. Now $F_x$ is a nonzero $\cal{O}_x$-module and
$\cal{O}_y$-flat, hence faithfully flat over~$\cal{O}_y$
by~\Ref{IV.2.2}. We can therefore apply~\Ref{IV.2.3}, which
completes the proof.

\begin{remarksstar}
As flatness is preserved by extension of the base, we see that
under the conditions of~\Ref{IV.6.5} $f$ is even
\emph{universally open}. I do not know however, when $Y$ is
integral and $X$ of finite type over~$Y$, whether $f$ induces on every
component $X_i$ of~$X$ an open morphism, or even only an
equidimensional one\footnote{The answer to the second question is
affirmative, the answer to the first negative even if $f$
is \'etale; \cf EGA~IV~12.1.1.5 and EGA~Err$_{\rm IV}$~33.},
\ie
% original p. 102
one all of whose fiber components have the same dimension (one knows
only that $X_i$ \emph{dominates} $Y$). The question is linked
to the following: let $A\to B$ be a local homomorphism of noetherian
local rings, such that $B$ is flat over~$A$ and~$\goth{m}B$
is an ideal of definition of~$B$ (which moreover implies
$\dim B=\dim A$). Is it true that for every minimal prime
ideal $\goth{p}_i$ of~$B$, one has $\dim B/\goth{p}_i=\dim B$?
Let us point out only that the answer to the first question
is negative when one replaces the flatness hypothesis
of~\Ref{IV.6.5} by the sole hypothesis that $f$ be
universally open.
\end{remarksstar}

\begin{lemma}
\label{IV.6.7}
Let $A$ be a noetherian integral domain, $B$ an
$A$-algebra of finite type, $M$ a $B$-module of finite type. Then
there exists a nonzero element $f$ of~$A$ such that $M_f$ is a
free (a fortiori flat) module over~$A_f$.
\end{lemma}

Let $K$ be the field of fractions of~$A$; then $B\otimes_A K$ is an
algebra of finite type over~$K$, and $M\otimes_A K$ a module of finite
type over the latter. Let $n$ be the dimension of the support of this
module; we shall reason by induction on~$n$.
If $n<0$,
\ie if $M\otimes_A K=0$, then taking a finite number of
generators of~$M$ over~$B$, we see that there exists an $f\in A$
which annihilates these generators, hence $M$, whence $M_f=0$ and we have
won. Suppose $n\ge 0$. We know that the $B$-module~$M$ admits
a composition series whose successive quotients are isomorphic
to modules $B/\goth{p}_i$, the $\goth{p}_i$ being
prime ideals of~$B$. As an extension of free modules is
free, we are reduced to the case where $M$ itself is of the
form $B/\goth{p}$, or again identical to $B$, $B$ being an
$A$-algebra that is an \emph{integral domain}. Applying the Noether
normalization lemma to the $K$-algebra $B\otimes_A K$, we
see easily that there exists a nonzero element $f$ of~$A$ such
that $B_f$ is integral over the subring $A_f[t_1,\dots,t_n]$, where
the $t_i$ are indeterminates. Hence we can already
suppose $B$ is integral over~$C=A[t_1,\dots,t_n]$; it is therefore a
torsion-free $C$-module of finite type. Let $m$ be its rank; there exists
therefore an exact sequence of~$C$-modules:
$$
0\to C^m\to B\to M'\to 0
$$
where $M'$ is a torsion $C$-module. It follows that the
Krull dimension of the~$C\otimes_A K$-module $M'\otimes_A K$ is
strictly less than the dimension~$n$ of~$C\otimes_A K$.
By the induction hypothesis, it follows that, on
condition of localizing with respect to a suitable nonzero $f$
of~$A$, we can suppose that $M'$ is a free
% original p. 103
$A$-module. On the other hand, $C^m$ is a free $A$-module. Hence $B$ is
then a free $A$-module; we are done.

\begin{lemma}
\label{IV.6.8}
Let $A$ be a noetherian ring, $B$ an algebra of finite type
over~$A$, $M$~a $B$-module of finite type, $\goth{p}$ a prime
ideal of~$B$, $\goth{q}$ the prime ideal that it induces on~$A$. We suppose $M_\goth{p}$ flat over~$A_{\goth{q}}$ (or over~$A$,
it is the same thing). Then there exists a $g\in B-\goth{p}$ such that
\begin{enumerate}
\item[(a)] $(M/\goth{q}M)_g$ is flat over~$A/\goth{q}$.
\item[(b)] $\Tor_1^A(M,A/\goth{q})_g=0$.
\end{enumerate}
\end{lemma}
Indeed, applying~\Ref{IV.6.7} to $(A/\goth(q), B/\goth{q}B,
M/\goth{q}M)$ we see first that there exists an~$f$ in $A-\goth{q}$
such that $(M/\goth{q}M)_f$ is flat over~$A/\goth{q}$. On the other hand,
as $M_\goth{p}$ is flat over~$A$, we have
$\Tor_1^A(M,A/\goth{q})_\goth{p}=\Tor_1^A(M_\goth{p},A/\goth{q})=0$,
hence, as $\Tor_1^A(M,A/\goth{q})$ is a $B$-module of finite type, there
exists a $g\in B-\goth{p}$ such that we have (b). We can then
(replacing $g$ by $gf$) suppose that we have at the same time (a),
which proves the corollary.

\begin{corollary}
\label{IV.6.9}
With the notation of~\Ref{IV.6.8}, for every prime ideal
$\goth{p}'$ of~$B$ containing~$\goth{p}$ and not containing~$g$,
$M_{\goth{p}'}$ is flat over~$A$ (or, what amounts to the same, over~$A_{\goth{q}'}$, where $\goth{q}'$ is the prime ideal of~$A$
induced by~$\goth{p}')$.
\end{corollary}
It suffices to apply criterion~\Ref{IV.5.6} (ii) to the system
$(A, B_{\goth{q}'}, \goth{q}, M_{\goth{q'}})$, using the
localization of the~$\Tor$.
\begin{theorem}
\label{IV.6.10}
Let $f\colon X\to Y$ be a morphism of finite type, with $Y$ locally
noetherian, and let $F$ be a coherent sheaf on~$X$. Let $U$
be the set of points $x\in X$ such that $F_x$ is flat over~$\cal{O}_{f(x)}$. Then $U$ is an \emph{open} set.
\end{theorem}

\subsubsection*{Proof} We can suppose $X$ and $Y$ affine, with rings $B$
and $A$, hence $F$ defined by a $B$-module $M$ of finite type. We
apply criterion~\Ref{IV.6.4}. Condition (a) is
verified trivially by \Ref{IV.1.2}~(i); it remains to
verify condition (b) of~\Ref{IV.6.4}. This is what was
done in lemma~\Ref{IV.6.8} and
corollary~\Ref{IV.6.9}.

In many questions, the following weaker form of
theorem~\Ref{IV.6.10} is sufficient (which already follows
from lemma~\Ref{IV.6.7}, and therefore requires neither the
technique
% original p. 104
of constructibles, nor theorem~\Ref{IV.5.6}):
\begin{corollary}
\label{IV.6.11}
Under the conditions of~\Ref{IV.6.10}, if we suppose $Y$ integral,
then there exists a nonempty open set~$V$ in $Y$ such that $F$ is flat
relative to $Y$ at all points of~$f^{-1}(V)$.
\end{corollary}
Indeed, the open set $U$ contains the fiber of the generic
point of~$Y$ (since the local ring of this point is a
field), hence it contains an open set of the form $f^{-1}(V)$, $X$
being of finite type over~$Y$. From~\Ref{IV.6.11}, one also
easily concludes the following result, where~$Y$ is assumed
noetherian (but not necessarily integral): there exists
a partition of~$Y$ into locally closed subsets $Y_i$
such that (equipping $Y_i$ with the induced reduced structure) $F$
induces on each $X_i=X\times_Y Y_i$ a sheaf flat with respect
to~$Y_i$.
